\section{An accepted three-fault logical error at $d=5$}
\label{app:d5-fault-v2}
The $d=5$ double-checking circuit also admits a small fault pattern
that is not caught by post-selection, leading to a logical error in the post-selected state. 
We take the ideal encoded input $\ket{T_L^\dagger}$ and place the three faults shown in
Fig.~\ref{fig:d5-fixed-faults}:
\begin{itemize}
\item $F_1$: $X_3$ after $\text{CNOT}_{3\to1}$ fold,
\item $F_2$: $X_{14}$ after $\text{CNOT}_{14\to15}$ in fold, and
\item $F_3$: $Y_{23}$ after $\text{CNOT}_{30\to23}$ in unfold.
\end{itemize}
No other location is faulty.
Each specified fault is one outcome of a two-qubit depolarising (e.g. $I\otimes X$, $X\otimes Y$ etc)
channel, with probability $p/15$ at its location.

The unfold propagates the first two faults to $X_0X_1X_3$ and
$X_9X_{14}X_{15}$. The ancillary factors $X_1$ and $X_{14}$ commute
with their $X$ readouts. The central and ancillary measurements
all return $+1$. Before the exit layer $D$, the data qubits pick up the error
\begin{equation}
 Q=X_0X_3X_9X_{15}Y_{23}.
 \label{eq:d5-fault-frame-v2}
\end{equation}
The exit gates expand $DQD^\dagger$ into 32 signed Pauli terms.
Only two survive the noiseless code projection $\Pi_0$:
\begin{equation}
 \Pi_0DQD^\dagger\Pi_0 =-\frac{X_L+Y_L}{\sqrt{32}}\Pi_0 =-\frac14N_L\Pi_0,
 \label{eq:d5-fault-projection-v2}
\end{equation}
where the logical operator $N_L=(X_L+Y_L)/\sqrt2$ acts on the nineteen-qubit code.
It anticommutes with $M_L=(X_L-Y_L)/\sqrt2$, so it maps the ideal
$M_L=+1$ magic state to the orthogonal $M_L=-1$ state.
The amplitude factor $1/4$ gives
\begin{equation}
 A_{\mathcal F}=B_{\mathcal F}=\frac1{16},
 \qquad \frac{B_{\mathcal F}}{A_{\mathcal F}}=1.
 \label{eq:d5-fixed-fault-probabilities-v2}
\end{equation}
Thus 1/16 of attempts with this fixed three-fault pattern
pass post-selection, and every accepted output exhibits a logical error.
This is an explicit weight-three failure contributing to the leading
term of Eq.~\eqref{eq:fullseries}. A ZX-calculus contraction of this weight-three fault pattern using techniques from~\cite{wan} further confirms this result.
Hence the $d=5$ cultivation circuit actually has fault distance $d_{\text{fault}} =3$ instead of $5$.

\begin{figure}[tp]\centering
\begin{adjustbox}{max width=.48\textwidth,max totalheight=.87\textheight}
\begingroup
\usetikzlibrary{decorations.pathreplacing}
\definecolor{faultorange}{RGB}{199,97,0}
\definecolor{faultblue}{RGB}{0,121,163}
\definecolor{faultpurple}{RGB}{123,55,156}
\begin{tikzpicture}[font=\scriptsize,line width=.45pt]
\begin{yquant}[register/minimum height=1.2mm,register/minimum depth=1.2mm,register/separation=.7mm,operator/minimum width=3.6mm,operator/separation=2.0mm,every wire/.style={draw=black!45}]
[name=row24] qubit {$24\,\ket{+}$} q24;
[name=row25] qubit {$25$} q25;
[name=row31] qubit {$31\,\ket{+}$} q31;
[name=row32] qubit {$32$} q32;
[name=row37] qubit {$37\,\ket{+}$} q37;
[name=row36] qubit {$36$} q36;
[name=row41] qubit {$41\,\ket{+}$} q41;
[name=row40] qubit {$40$} q40;
[name=row26] qubit {$26\,\ket{+}$} q26;
[name=row19] qubit {$19$} q19;
[name=row33] qubit {$33\,\ket{+}$} q33;
[name=row34] qubit {$34$} q34;
[name=row39] qubit {$39\,\ket{+}$} q39;
[name=row38] qubit {$38$} q38;
[name=row20] qubit {$20\,\ket{+}$} q20;
[name=row27] qubit {$27$} q27;
[name=row28] qubit {$28\,\ket{+}$} q28;
[name=row21] qubit {$21$} q21;
[name=row22] qubit {$22\,\ket{+}$} q22;
[name=row29] qubit {$29$} q29;
[name=row30] qubit {$30\,\ket{+}$} q30;
[name=row23] qubit {$23$} q23;
[name=row18] qubit {$18\,\ket{+}$} q18;
[name=row17] qubit {$17$} q17;
[name=row12] qubit {$12\,\ket{+}$} q12;
[name=row11] qubit {$11$} q11;
[name=row8] qubit {$8\,\ket{+}$} q8;
[name=row13] qubit {$13$} q13;
[name=row14] qubit {$14\,\ket{+}$} q14;
[name=row9] qubit {$9$} q9;
[name=row10] qubit {$10\,\ket{+}$} q10;
[name=row15] qubit {$15$} q15;
[name=row6] qubit {$6\,\ket{+}$} q6;
[name=row7] qubit {$7$} q7;
[name=row2] qubit {$2\,\ket{+}$} q2;
[name=row3] qubit {$3$} q3;
[name=row1] qubit {$1\,\ket{+}$} q1;
[name=row0] qubit {$0$} q0;
[name=entry] box {$T$} q0;
box {$T$} q11;
box {$T$} q17;
box {$T$} q19;
box {$T$} q36;
box {$T$} q40;
box {$T^\dagger$} q3;
box {$T^\dagger$} q7;
box {$T^\dagger$} q9;
box {$T^\dagger$} q13;
box {$T^\dagger$} q15;
box {$T^\dagger$} q21;
box {$T^\dagger$} q23;
box {$T^\dagger$} q25;
box {$T^\dagger$} q27;
box {$T^\dagger$} q29;
box {$T^\dagger$} q32;
box {$T^\dagger$} q34;
box {$T^\dagger$} q38;
align q24,q25,q31,q32,q37,q36,q41,q40,q26,q19,q33,q34,q39,q38,q20,q27,q28,q21,q22,q29,q30,q23,q18,q17,q12,q11,q8,q13,q14,q9,q10,q15,q6,q7,q2,q3,q1,q0;
[name=g394n0] cnot q0 | q1;
[name=g394n1] cnot q3 | q2;
[name=g394n2] cnot q7 | q6;
[name=g394n3] cnot q13 | q8;
[name=g394n4] cnot q15 | q10;
[name=g394n5] cnot q11 | q12;
[name=g394n6] cnot q9 | q14;
[name=g394n7] cnot q17 | q18;
[name=g394n8] cnot q27 | q20;
[name=g394n9] cnot q29 | q22;
[name=g394n10] cnot q25 | q24;
[name=g394n11] cnot q19 | q26;
[name=g394n12] cnot q21 | q28;
[name=g394n13] cnot q23 | q30;
[name=g394n14] cnot q32 | q31;
[name=g394n15] cnot q34 | q33;
[name=g394n16] cnot q36 | q37;
[name=g394n17] cnot q38 | q39;
[name=g394n18] cnot q40 | q41;
align q24,q25,q31,q32,q37,q36,q41,q40,q26,q19,q33,q34,q39,q38,q20,q27,q28,q21,q22,q29,q30,q23,q18,q17,q12,q11,q8,q13,q14,q9,q10,q15,q6,q7,q2,q3,q1,q0;
[name=g397n0] cnot q1 | q3;
[name=g397n1] cnot q15 | q14;
[name=g397n2] cnot q30 | q29;
[name=g397n3] cnot q38 | q33;
[name=g397n4] cnot q41 | q37;
align q24,q25,q31,q32,q37,q36,q41,q40,q26,q19,q33,q34,q39,q38,q20,q27,q28,q21,q22,q29,q30,q23,q18,q17,q12,q11,q8,q13,q14,q9,q10,q15,q6,q7,q2,q3,q1,q0;
[name=F1,draw=faultorange,fill=faultorange!13,text=faultorange,line width=.8pt] box {$X_{F_1}$} q3;
setstyle {draw=faultorange,line width=1.1pt} q3;
[name=F2,draw=faultblue,fill=faultblue!13,text=faultblue,line width=.8pt] box {$X_{F_2}$} q14;
setstyle {draw=faultblue,line width=1.1pt} q14;
align q24,q25,q31,q32,q37,q36,q41,q40,q26,q19,q33,q34,q39,q38,q20,q27,q28,q21,q22,q29,q30,q23,q18,q17,q12,q11,q8,q13,q14,q9,q10,q15,q6,q7,q2,q3,q1,q0;
align q24,q25,q31,q32,q37,q36,q41,q40,q26,q19,q33,q34,q39,q38,q20,q27,q28,q21,q22,q29,q30,q23,q18,q17,q12,q11,q8,q13,q14,q9,q10,q15,q6,q7,q2,q3,q1,q0;
[name=g400n0] cnot q3 | q7;
[name=g400n1] cnot q14 | q13;
[name=g400n2] cnot q33 | q26;
[name=g400n3] cnot q29 | q28;
[name=g400n4] cnot q37 | q32;
align q24,q25,q31,q32,q37,q36,q41,q40,q26,q19,q33,q34,q39,q38,q20,q27,q28,q21,q22,q29,q30,q23,q18,q17,q12,q11,q8,q13,q14,q9,q10,q15,q6,q7,q2,q3,q1,q0;
[name=g403n0] cnot q13 | q12;
[name=g403n1] cnot q32 | q25;
[name=g403n2] cnot q28 | q27;
align q24,q25,q31,q32,q37,q36,q41,q40,q26,q19,q33,q34,q39,q38,q20,q27,q28,q21,q22,q29,q30,q23,q18,q17,q12,q11,q8,q13,q14,q9,q10,q15,q6,q7,q2,q3,q1,q0;
[name=g406n0] cnot q7 | q12;
[name=g406n1] cnot q27 | q26;
align q24,q25,q31,q32,q37,q36,q41,q40,q26,q19,q33,q34,q39,q38,q20,q27,q28,q21,q22,q29,q30,q23,q18,q17,q12,q11,q8,q13,q14,q9,q10,q15,q6,q7,q2,q3,q1,q0;
[name=g409n0] cnot q12 | q18;
[name=g409n1] cnot q26 | q25;
align q24,q25,q31,q32,q37,q36,q41,q40,q26,q19,q33,q34,q39,q38,q20,q27,q28,q21,q22,q29,q30,q23,q18,q17,q12,q11,q8,q13,q14,q9,q10,q15,q6,q7,q2,q3,q1,q0;
[name=g412n0] cnot q18 | q25;
align q24,q25,q31,q32,q37,q36,q41,q40,q26,q19,q33,q34,q39,q38,q20,q27,q28,q21,q22,q29,q30,q23,q18,q17,q12,q11,q8,q13,q14,q9,q10,q15,q6,q7,q2,q3,q1,q0;
[name=hubmeasure,type=qubit] measure {$X$} q25;
[name=hubreset] box {$R_X$} q25;
align q24,q25,q31,q32,q37,q36,q41,q40,q26,q19,q33,q34,q39,q38,q20,q27,q28,q21,q22,q29,q30,q23,q18,q17,q12,q11,q8,q13,q14,q9,q10,q15,q6,q7,q2,q3,q1,q0;
align q24,q25,q31,q32,q37,q36,q41,q40,q26,q19,q33,q34,q39,q38,q20,q27,q28,q21,q22,q29,q30,q23,q18,q17,q12,q11,q8,q13,q14,q9,q10,q15,q6,q7,q2,q3,q1,q0;
[name=g420n0] cnot q18 | q25;
align q24,q25,q31,q32,q37,q36,q41,q40,q26,q19,q33,q34,q39,q38,q20,q27,q28,q21,q22,q29,q30,q23,q18,q17,q12,q11,q8,q13,q14,q9,q10,q15,q6,q7,q2,q3,q1,q0;
[name=g423n0] cnot q12 | q18;
[name=g423n1] cnot q26 | q25;
align q24,q25,q31,q32,q37,q36,q41,q40,q26,q19,q33,q34,q39,q38,q20,q27,q28,q21,q22,q29,q30,q23,q18,q17,q12,q11,q8,q13,q14,q9,q10,q15,q6,q7,q2,q3,q1,q0;
[name=g426n0] cnot q7 | q12;
[name=g426n1] cnot q27 | q26;
align q24,q25,q31,q32,q37,q36,q41,q40,q26,q19,q33,q34,q39,q38,q20,q27,q28,q21,q22,q29,q30,q23,q18,q17,q12,q11,q8,q13,q14,q9,q10,q15,q6,q7,q2,q3,q1,q0;
[name=g429n0] cnot q13 | q12;
[name=g429n1] cnot q32 | q25;
[name=g429n2] cnot q28 | q27;
align q24,q25,q31,q32,q37,q36,q41,q40,q26,q19,q33,q34,q39,q38,q20,q27,q28,q21,q22,q29,q30,q23,q18,q17,q12,q11,q8,q13,q14,q9,q10,q15,q6,q7,q2,q3,q1,q0;
[name=g432n0] cnot q3 | q7;
[name=g432n1] cnot q14 | q13;
[name=g432n2] cnot q33 | q26;
[name=g432n3] cnot q29 | q28;
[name=g432n4] cnot q37 | q32;
align q24,q25,q31,q32,q37,q36,q41,q40,q26,q19,q33,q34,q39,q38,q20,q27,q28,q21,q22,q29,q30,q23,q18,q17,q12,q11,q8,q13,q14,q9,q10,q15,q6,q7,q2,q3,q1,q0;
[name=g435n0,style={color=faultorange,draw=faultorange,line width=1.1pt},control style={fill=faultorange}] cnot q1 | q3;
setstyle {draw=faultorange,line width=1.1pt} q1;
[name=g435n1,style={color=faultblue,draw=faultblue,line width=1.1pt},control style={fill=faultblue}] cnot q15 | q14;
setstyle {draw=faultblue,line width=1.1pt} q15;
[name=g435n2] cnot q30 | q29;
[name=g435n3] cnot q38 | q33;
[name=g435n4] cnot q41 | q37;
align q24,q25,q31,q32,q37,q36,q41,q40,q26,q19,q33,q34,q39,q38,q20,q27,q28,q21,q22,q29,q30,q23,q18,q17,q12,q11,q8,q13,q14,q9,q10,q15,q6,q7,q2,q3,q1,q0;
[name=g438n0,style={color=faultorange,draw=faultorange,line width=1.1pt},control style={fill=faultorange}] cnot q0 | q1;
setstyle {draw=faultorange,line width=1.1pt} q0;
[name=g438n1] cnot q3 | q2;
[name=g438n2] cnot q7 | q6;
[name=g438n3] cnot q13 | q8;
[name=g438n4] cnot q15 | q10;
[name=g438n5] cnot q11 | q12;
[name=g438n6,style={color=faultblue,draw=faultblue,line width=1.1pt},control style={fill=faultblue}] cnot q9 | q14;
setstyle {draw=faultblue,line width=1.1pt} q9;
[name=g438n7] cnot q17 | q18;
[name=g438n8] cnot q27 | q20;
[name=g438n9] cnot q29 | q22;
[name=g438n10] cnot q25 | q24;
[name=g438n11] cnot q19 | q26;
[name=g438n12] cnot q21 | q28;
[name=g438n13] cnot q23 | q30;
[name=g438n14] cnot q32 | q31;
[name=g438n15] cnot q34 | q33;
[name=g438n16] cnot q36 | q37;
[name=g438n17] cnot q38 | q39;
[name=g438n18] cnot q40 | q41;
align q24,q25,q31,q32,q37,q36,q41,q40,q26,q19,q33,q34,q39,q38,q20,q27,q28,q21,q22,q29,q30,q23,q18,q17,q12,q11,q8,q13,q14,q9,q10,q15,q6,q7,q2,q3,q1,q0;
[name=F3,draw=faultpurple,fill=faultpurple!13,text=faultpurple,line width=.8pt] box {$Y_{F_3}$} q23;
setstyle {draw=faultpurple,line width=1.1pt} q23;
align q24,q25,q31,q32,q37,q36,q41,q40,q26,q19,q33,q34,q39,q38,q20,q27,q28,q21,q22,q29,q30,q23,q18,q17,q12,q11,q8,q13,q14,q9,q10,q15,q6,q7,q2,q3,q1,q0;
[name=exit] box {$T^\dagger$} q0;
setstyle {draw=black!45,line width=.45pt} q0;
box {$T^\dagger$} q11;
box {$T^\dagger$} q17;
box {$T^\dagger$} q19;
box {$T^\dagger$} q36;
box {$T^\dagger$} q40;
box {$T$} q3;
setstyle {draw=black!45,line width=.45pt} q3;
box {$T$} q7;
box {$T$} q9;
setstyle {draw=black!45,line width=.45pt} q9;
box {$T$} q13;
box {$T$} q15;
setstyle {draw=black!45,line width=.45pt} q15;
box {$T$} q21;
box {$T$} q23;
setstyle {draw=black!45,line width=.45pt} q23;
box {$T$} q25;
box {$T$} q27;
box {$T$} q29;
box {$T$} q32;
box {$T$} q34;
box {$T$} q38;
align q24,q25,q31,q32,q37,q36,q41,q40,q26,q19,q33,q34,q39,q38,q20,q27,q28,q21,q22,q29,q30,q23,q18,q17,q12,q11,q8,q13,q14,q9,q10,q15,q6,q7,q2,q3,q1,q0;
[type=qubit] measure {$X$} q20;
[type=qubit] measure {$X$} q22;
[type=qubit] measure {$X$} q8;
[type=qubit] measure {$X$} q10;
[type=qubit] measure {$X$} q2;
[type=qubit] measure {$X$} q6;
[type=qubit] measure {$X$} q24;
[type=qubit] measure {$X$} q31;
[type=qubit] measure {$X$} q39;
[type=qubit] measure {$X$} q1;
setstyle {draw=black!45,line width=.45pt} q1;
[type=qubit] measure {$X$} q12;
[type=qubit] measure {$X$} q18;
[type=qubit] measure {$X$} q37;
[type=qubit] measure {$X$} q33;
[type=qubit] measure {$X$} q14;
setstyle {draw=black!45,line width=.45pt} q14;
[type=qubit] measure {$X$} q26;
[type=qubit] measure {$X$} q28;
[type=qubit] measure {$X$} q30;
[type=qubit] measure {$X$} q41;
\end{yquant}
\coordinate (heading) at ($(row24)+(0,8mm)$);
\node[anchor=south,font=\small] at (entry |- heading) {$D^\dagger$};
\node[anchor=south,font=\small] at (exit |- heading) {$D$};
\draw[black!60,decorate,decoration={brace,amplitude=2.5pt}] (g394n0 |- heading) -- (g412n0 |- heading) node[midway,above=3pt,font=\small] {fold};
\draw[black!60,decorate,decoration={brace,amplitude=2.5pt}] (g420n0 |- heading) -- (F3.east |- heading) node[midway,above=3pt,font=\small] {unfold};
\end{tikzpicture}

\endgroup

\end{adjustbox}
\caption{Three faults in the $d=5$ double-checking circuit.
Coloured boxes mark the faults; coloured wires trace their propagation.
Conditioned on this pattern, $A_{\mathcal F}=B_{\mathcal F}=1/16$.}
\label{fig:d5-fixed-faults}
\end{figure}
\FloatBarrier

